\chapter{Statistical Methodology and Replication Analysis}
\label{app:statistical_methods}

This appendix presents the complete mathematical formulations, hypotheses tests, and computational replication procedures underpinning both the historical human evaluator benchmark by Allodi et al.~\cite{santos2017estimating} and the HelpUS validation engine.

\section{Hypothesis Formulation and Fisher's Exact Test}
\label{sec:fisher_formulation}
To test whether human evaluators systematically fail when assessing vulnerabilities in complex environments, the evaluation is formalized as a series of $2 \times 2$ contingency tables. Let the assessment outcome for any participant evaluating vulnerability $v$ in network context $cxty \in \{\text{Flat}, \text{Segmented}\}$ be a binary random variable $Y \in \{0, 1\}$, where $Y = 1$ denotes agreement with the verified expert standard (correct) and $Y = 0$ denotes disagreement (incorrect).

\subsection{Network Complexity Hypothesis ($H2$)}
The null hypothesis states that the probability of correct assessment is independent of network topology:
\begin{align}
    H2_0: \quad & P(Y = 1 \mid \text{Flat}) = P(Y = 1 \mid \text{Segmented}) \\
    H2_1: \quad & P(Y = 1 \mid \text{Flat}) \neq P(Y = 1 \mid \text{Segmented})
\end{align}

Table~\ref{tab:app_contingency_network} summarizes the empirical contingency counts observed across the $n=29$ evaluator cohort ($N=348$ total scored instances).

\begin{table}[htbp]
\centering
\caption{Contingency Table of Assessment Correctness by Network Complexity}
\label{tab:app_contingency_network}
\begin{tabular}{lccc}
\toprule
\textbf{Network Architecture} & \textbf{Correct ($Y=1$)} & \textbf{Incorrect ($Y=0$)} & \textbf{Total} \\
\midrule
Flat (Simple)         & $a = 96$  & $b = 20$  & $n_1 = 116$ \\
Segmented (Complex)   & $c = 136$ & $d = 96$  & $n_2 = 232$ \\
\midrule
\textbf{Total}        & $m_1 = 232$ & $m_2 = 116$ & $N = 348$ \\
\bottomrule
\end{tabular}
\end{table}

Fisher's exact test computes the hypergeometric probability of observing the table marginals under $H2_0$:
\begin{equation}
    p = \frac{\binom{a+b}{a} \binom{c+d}{c}}{\binom{N}{a+c}} = \frac{\binom{116}{96} \binom{232}{136}}{\binom{348}{232}}
\end{equation}
The reproducibility script bundled with the project computes the two-sided Fisher exact test and both sample and conditional odds-ratio quantities. The generated values are:
\begin{align}
    p_{\text{Fisher, two-sided}} &= \HistFisherPScientific, \\n    \text{OR}_{\text{sample}} &= \HistSampleOR, \\n    \text{OR}_{\text{conditional MLE}} &= \HistConditionalOR.
\end{align}
Using the conditional exact method implemented by the pinned statistical toolchain, the $95\%$ confidence interval is:
\begin{equation}
    \text{CI}_{95\%,\text{conditional exact}} = [\HistConditionalCILow, \HistConditionalCIHigh].
\end{equation}
The method label is reported explicitly because this interval is not the same object as the simple asymptotic Wald interval. The historical null hypothesis of independence is rejected for this contingency table.

\section{Permutation Two-Way ANOVA for Co-Factor Invariance}
\label{sec:permutation_anova}
To ensure that the human error rate was not driven by individual demographic confounding factors (e.g., years of prior industry experience or self-reported security proficiency), the evaluation executed a non-parametric \textbf{Permutation Two-Way ANOVA}.

Let $F_{\text{obs}}$ denote the observed test statistic for the interaction between network scenario complexity and participant security experience. The null distribution is generated by randomly permuting the participant labels across $B = 10,000$ resamples:
\begin{equation}
    p_{\text{perm}} = \frac{1}{B} \sum_{b=1}^{B} \mathbb{I}\left( F^{(b)} \ge F_{\text{obs}} \right)
\end{equation}
The original participant-level data and complete permutation outputs are required to reproduce this co-factor analysis. Until those artifacts are included, this result should be treated as a historical-study claim rather than a newly reproduced HelpUS result.

\section{Monte Carlo Dispersal Simulation Algorithm}
\label{sec:monte_carlo_algo}
The synthetic score-space dispersal is generated by Algorithm~\ref{alg:monte_carlo}. It is a sensitivity simulation and is not the empirical distribution of human participant scores.

\begin{algorithm}[H]
\caption{Monte Carlo Scoring Dispersal Simulation}
\label{alg:monte_carlo}
\begin{algorithmic}[1]
\Require Base Score $S_{\text{base}}$, Permutation Count $N = 10,000$, Topology Type $\mathcal{T}$
\Ensure Score Distribution Vector $\mathbf{S}_{\text{dist}}$
\State $\mathbf{S}_{\text{dist}} \gets \emptyset$
\For{$i = 1$ to $N$}
    \State Sample $\text{CR} \sim \mathcal{U}\{\text{High}, \text{Medium}, \text{Low}\}$
    \State Sample $\text{IR} \sim \mathcal{U}\{\text{High}, \text{Medium}, \text{Low}\}$
    \State Sample $\text{AR} \sim \mathcal{U}\{\text{High}, \text{Medium}, \text{Low}\}$
    \If{$\mathcal{T} = \text{Segmented}$}
        \State Sample $\text{MAV} \sim \mathcal{U}\{\text{Network}, \text{Adjacent}, \text{Local}\}$
    \Else
        \State $\text{MAV} \gets \text{Network}$
    \EndIf
    \State $S_{\text{sim}} \gets \text{ComputeCVSS40}(S_{\text{base}}, \text{MAV}, \text{CR}, \text{IR}, \text{AR})$
    \State $\mathbf{S}_{\text{dist}} \gets \mathbf{S}_{\text{dist}} \cup \{S_{\text{sim}}\}$
\EndFor
\State \Return $\mathbf{S}_{\text{dist}}$
\end{algorithmic}
\end{algorithm}
